%% kapa: exam=1 round=37 year=2026 period=1 form=A subject=real_estate subject_ko=부동산학원론 seq=22 number=102
%% source: Q-Net 감정평가사 기출문제 (cst00309) artlSeq=5258924 file=raw/1008/5258924/2026년 제37회 감정평가사 1차 1교시.pdf
%% conversion: kordoc 4.24.51 -> tools/gen1.py (auto); review=proofread-sonnet; flags=none
\begin{question}{102}
위험회피형 투자자가 평균분산지배원리를 적용하여 투자안을 검토하고 있다. 이 투자자의 효용함수는 2차 함수 형태이며, 각 자산의 수익률은 정규분포를 따른다. 각 자산의 기대수익률과 표준편차가 다음과 같을 때 옳은 것은?
\par
{\small\centering
\begin{tabular}{|c|c|c|c|}\hline
구분 & 자산A & 자산B & 자산C \\ \hline
기대수익률 & 8 \% & 12 \% & ? \\ \hline
표준편차 & 4 \% & ? & 10 \% \\ \hline
\end{tabular}\par}
\begin{choices}
  \item B의 표준편차가 2 \%이면, A가 B를 지배한다.
  \item C의 기대수익률이 12 \%이면, C가 A를 지배한다.
  \item B의 표준편차가 4 \%이면, A와 B 사이에는 지배관계가 성립하지 않는다.
  \item B의 표준편차가 10 \%이고 C의 기대수익률이 16 \%이면, B와 C 사이에는 지배관계가 성립하지 않는다.
  \item B의 표준편차가 8 \%이고 C의 기대수익률이 16 \%이면, A, B, C 어느 자산도 다른 자산을 지배하지 않는다.
\end{choices}
\end{question}
\answer{5}
